% Short opener for the Evaluation section.

We evaluate \biotest{} along three complementary axes.
\S\ref{sec:coverage} reports final line coverage of the corpus
produced by \biotest{} relative to corpora produced by mainstream
coverage-guided fuzzers. \S\ref{sec:mutation-adequacy} grades the
same corpus on mutation adequacy under matched mutation engines.
\S\ref{sec:bug-bench} reports a real-bug benchmark of $32$
verified historical parser bugs. Mutation kills correlate with
real-fault detection at
$r{\approx}0.7$~\citep{just2014mutants}, so the three axes
together provide construct triangulation.

The evaluation spans six $(\mathrm{SUT, format})$ pairs:
htsjdk/VCF, htsjdk/SAM (Java), vcfpy/VCF, biopython/SAM (Python),
noodles/VCF (Rust), and seqan3/SAM (C++). All three axes use the
same matched protocol: identical mutation engine, target-class
scope, kill predicate, and wall-clock budget per pair, so only
the corpus differs between \biotest{} and each baseline. Every
baseline corpus is trimmed to $\le 200$ files by an
AFL-cmin-style minimiser, using per-file kill-cardinality on
Python SUTs~\citep{vikram2023guiding} and an SUT-agnostic
outcome-fingerprint elsewhere. The trim adds $+1$ to $+2$\,pp on
pairs where it actively engages. As an uninformed-fuzzing
reference~\citep{klees2018fuzz} we add a Pure-Random byte emitter
under identical conditions and report it as a per-pair floor.
All runs execute inside the reproducibility container described
in \S\ref{sec:implementation}.
